% Teil II: der durchgerechnete Weg.

% Gedankenschritt: Warum der Quotient die entscheidende Probe ist und wie
% man die Rückrechnung als Kontrolle einsetzt.

\subsection*{Gegeben und gesucht}
Zwei Grössen mit den Wertepaaren

\begin{center}
$X$: $4,\ 8,\ 10,\ 12$ \quad und \quad $Y$: $6,\ 12,\ 15,\ 18$.
\end{center}

Gesucht: Ist $Y$ direkt proportional zu $X$? Gegebenenfalls die
Proportionalitätskonstante angeben.

\subsection*{Ansatz und der entscheidende Gedankenschritt}
Direkte Proportionalität bedeutet: Es gibt eine Zahl $k$ mit $Y=kX$ für
alle Paare. Die Frage ``ist das so?`` übersetzt sich also in die
Rechenprobe ``sind alle Quotienten $k_i=Y_i/X_i$ gleich?``

\textbf{Warum der Quotient und nicht der Unterschied?} Zwei Grössen mit
konstantem \emph{Unterschied} (z.\,B. $Y=X+3$) wären additiv verknüpft,
nicht multiplikativ. Proportionalität ist eine Frage nach einem
\emph{Faktor}, und den extrahiert genau die Division. Der Quotient prüft
also gezielt die Vermutung ``Mal-Faktor statt Plus-Konstante``.

\subsection*{Rechnung}
\begin{align}
  \frac{6}{4}&=\frac32,\qquad \frac{12}{8}=\frac32,\qquad
  \frac{15}{10}=\frac32,\qquad \frac{18}{12}=\frac32 .
\end{align}
Alle vier Quotienten sind gleich $k=\tfrac32$. (Beachte: Das war vorab
nicht sicher --- z.\,B. ist $10/4=\tfrac52$, aber wir testen ja
$Y/X=15/10$, nicht $X/Y$.)

\begin{ergebnis}
$Y$ ist direkt proportional zu $X$ mit der Proportionalitätskonstante
$k=\dfrac32$, d.h. $Y=\tfrac32\,X$.
\end{ergebnis}

\begin{kontrolle}
\textbf{Rückrechnung (eine starke Probe):} $k\cdot X$: $\tfrac32\cdot4=6$,
$\tfrac32\cdot8=12$, $\tfrac32\cdot10=15$, $\tfrac32\cdot12=18$ -- alle
gegebenen $Y$-Werte werden reproduziert. Damit ist die Aussage nicht nur
``plausibel``, sondern \emph{bewiesen} an allen Daten.

\textbf{Grenzfall (was wäre, wenn $\neq$):} Hätte eines der Paare, etwa
$(10,16)$, den Quotienten $16/10=8/5\neq3/2$, so wäre die direkte
Proportionalität \emph{sofort} zu verwerfen --- ein Gegenpaar genügt.
Genau dieser ``minimale Widerlegungstest`` steckt hinter der Regel ``alle
Paare prüfen``. Auch der Fall $X=0$ mit $Y\neq0$ würde die Aussage
unmöglich machen (keine Ursprungsgerade).
\end{kontrolle}

\begin{merksatz}
Proportionalität erkennt man am konstanten Quotienten $Y/X$;
der gemeinsame Quotient ist die Proportionalitätskonstante. Die
Rückrechnung $kX=Y$ an jedem Paar ist der zuverlässige Abschluss.
\end{merksatz}
